\subsection{Sự chuyển dịch trọng tâm trong SDLC}

Sự phát triển của Generative AI đang làm thay đổi cách thức thực hiện nhiều hoạt động trong vòng đời phát triển phần mềm (SDLC). Trước đây, kỹ sư phần mềm dành phần lớn thời gian trực tiếp thực hiện các tác vụ như viết mã nguồn, tìm lỗi, kiểm thử và xử lý các vấn đề kỹ thuật. Khi các công cụ AI có khả năng hỗ trợ tạo mã, giải thích mã, phát hiện lỗi và đề xuất giải pháp ngày càng phổ biến, một phần công việc thực thi có thể được AI hỗ trợ hoặc tự động hóa.

Tuy nhiên, điều này không đồng nghĩa với việc vai trò của kỹ sư phần mềm bị thay thế hoàn toàn. Trọng tâm công việc có xu hướng chuyển từ trực tiếp thực hiện từng tác vụ kỹ thuật sang xác định vấn đề, phân rã yêu cầu, hướng dẫn AI, đánh giá kết quả, kiểm thử/hiệu chỉnh, tích hợp và chịu trách nhiệm đối với sản phẩm cuối cùng. Kỹ sư cần có khả năng xác định khi nào nên sử dụng AI, cung cấp ngữ cảnh và yêu cầu phù hợp, đồng thời kiểm tra xem kết quả do AI tạo ra có đáp ứng đúng yêu cầu hay không.

Sự chuyển dịch này khiến các kỹ năng truyền thống vẫn giữ vai trò quan trọng. Kiến thức về lập trình, kiến trúc phần mềm, kiểm thử, bảo mật và quản lý yêu cầu giúp kỹ sư đánh giá được chất lượng của kết quả do AI tạo ra thay vì chỉ chấp nhận kết quả một cách máy móc. Đồng thời, các kỹ năng mới như làm việc với AI, phân rã vấn đề, kiểm chứng đầu ra và đánh giá rủi ro trở nên quan trọng hơn trong quy trình phát triển.

Có thể khái quát sự thay đổi này như sau:

\begin{itemize}
    \item \textbf{Workflow truyền thống:} Yêu cầu $\rightarrow$ Phân tích $\rightarrow$ Thiết kế $\rightarrow$ Lập trình $\rightarrow$ Kiểm thử $\rightarrow$ Sửa lỗi $\rightarrow$ Hoàn thiện.
    \item \textbf{Workflow có AI hỗ trợ:} Yêu cầu $\rightarrow$ Phân rã vấn đề $\rightarrow$ Giao nhiệm vụ cho AI $\rightarrow$ Đánh giá đầu ra $\rightarrow$ Kiểm thử/hiệu chỉnh $\rightarrow$ Tích hợp và chịu trách nhiệm.
\end{itemize}

Như vậy, AI chủ yếu làm thay đổi cách kỹ sư thực hiện công việc, trong khi con người vẫn giữ vai trò quan trọng trong việc xác định mục tiêu, đưa ra quyết định và kiểm soát chất lượng.
